% `tex.web` §§336, 392, and 396: a macro call is abandoned when its
% compulsory parameter text does not match, when an outer control sequence is
% found at any brace depth, or when a non-long argument contains `\par`.
% Captured arguments are discarded and the replacement text is never entered.
\catcode`\{=1 \catcode`\}=2 \catcode`\#=6
\tracingonline=1 \tracingcommands=2 \tracingmacros=1
\outer\def\stop{}
\def\need A#1{\message{BAD:need}}
\need B
\need\stop
\def\undelimited#1{\message{BAD:undelimited}}
\undelimited{x\stop}
\def\delimited#1X{\message{BAD:delimited}}
\delimited a\stop X
% §400 traces each completed argument immediately. Thus the first argument is
% visible before an outer token makes the second one run away.
\def\ordered#1#2X{\message{BAD:ordered}}
\ordered A\stop X
% `\meaning` uses the current escape character for macro prefixes. LaTeX 2e's
% command-hook builder parses exactly this string with `/` as a delimiter.
\escapechar=`/
\long\def\prefixprobe#1{}
\message{PROBE:prefix=\meaning\prefixprobe}
\escapechar=`\\
% §396 compares the literal `\par` token, not all aliases of its meaning.
\let\paralias=\par
\def\aliasarg#1{\message{PROBE:par-alias-accepted}}
\aliasarg\paralias
\message{PROBE:matching-recovered}
\end
